\documentclass[11pt,a4paper]{article}
\usepackage{cmap}
\usepackage[utf8]{inputenc}
\usepackage[T1]{fontenc}
\usepackage{lmodern}
\usepackage{microtype}
\usepackage[margin=1.0in]{geometry}
\usepackage{amsmath,amssymb,amsthm,mathtools}
\usepackage{booktabs,array,tabularx}
\usepackage[svgnames]{xcolor}
\usepackage[obeyspaces,spaces]{url}
\usepackage[colorlinks=true,linkcolor=DarkBlue,citecolor=DarkBlue,urlcolor=DarkBlue]{hyperref}

\urlstyle{tt}
\def\UrlFont{\ttfamily\upshape}
\renewcommand{\arraystretch}{1.3}
\setlength{\tabcolsep}{5pt}
\sloppy
\setlength{\emergencystretch}{4em}
\hfuzz=4.0pt
\vfuzz=4.0pt

\newtheorem{theorem}{Theorem}[section]
\theoremstyle{definition}
\newtheorem{definition}[theorem]{Definition}
\theoremstyle{remark}
\newtheorem{remark}{Remark}[section]

\title{\textbf{Behavioural Profiles Are Not Heritable Sequences:\\
a Refuted Transfer of Phylogenetic Recombination Detection\\
to Genetic Programming}}
\author{Xavier Callens\\
\small SocrateAI Scientific Research \& Formal Verification Laboratory}
\date{Edition 1, 24 October 2026. \textbf{Not peer reviewed.}}

\begin{document}
\maketitle

\begin{abstract}
Evolutionary program synthesis in the AlphaEvolve family applies mutation and
selection to populations of programs, usually with a crossover operator whose
contribution is asserted rather than measured. Phylogenetics has an exact
instrument for precisely that question: a quartet certificate refutes the
four-point condition on four taxa with an integer comparison a proof kernel can
recheck, and a recombinant genealogy violates that condition where a clonal one
does not. This paper proposed transferring the instrument, tested the transfer,
and reports its failure together with the reason.

The transfer looked clean. A candidate program's \emph{behavioural profile} ---
its outcome at each of $T$ probes, drawn from \{correct, incorrect, raised,
non-finite\} with a gap symbol for an inadmissible probe --- is a four-state
sequence with missing data, which is exactly the model a nucleotide alignment
assumes. A $\text{programs} \times \text{probes}$ matrix is therefore an
alignment, and the certificate machinery runs on it unmodified. A natural null
comes for free: a population varied only by mutation from single parents is a
tree \emph{by construction}.

On synthetic profiles the transfer appeared to work, separating clonal from
recombinant arms by margin across a usable window of mutation rates. On real
program populations --- expression trees over a closed grammar, evolved against a
rational target, every individual executed at 600 probes --- it fails. Three
pre-registered controls fail in the same direction: the recombinant arm exceeds
the clonal floor in $0$ of $24$ runs, $0$ of $23$ when restricted to comparable
best fitness, and $0$ of $22$ when restricted to comparable \emph{realised
behavioural spread}. Stratifying all $49$ populations by spread then explains
everything: above a spread of $0.05$, $4$ of $4$ clonal populations carry
certificates and so does $1$ of $1$ recombinant; below $0.01$, none of $29$ does
in either arm. Behavioural spread accounts for certificate presence and arm
membership adds nothing.

The diagnosis is specific and it indicts the proposal rather than the
implementation. The four-point condition detects non-tree-like structure in a
distance matrix; it detects \emph{recombination} only when the characters are
independently heritable units, because recombination's signature is mosaic
ancestry \emph{along} the sequence. A probe outcome is not inherited: it is a
consequence of the whole program, so two unrelated point mutations routinely
alter overlapping sets of probe outcomes and manufacture the same pattern. The
synthetic validation concealed this because its profiles were generated by
mutating profile \emph{strings}, which builds site-wise heritability in by
construction --- it confirmed that the method detects recombination in
site-wise-heritable data, which was never in doubt.

We report what survives: the proof-carrying fitness ladder, which is independent
of this result; the twelve measurement guards, three of which caught this
failure; and a released dataset of $49$ evolved program populations with
behavioural profiles, certificates and fitness traces.
\end{abstract}

\section{What was proposed, and why it was plausible}

Systems in the AlphaEvolve family evolve programs with an LLM or with
syntax-level operators, score them against an evaluator, and keep what scores
better. Most carry a crossover operator. Whether crossover contributes is an old
and unsettled question in genetic programming, and in practice it is answered by
assertion: the operator is present, so recombination is assumed to occur.

Phylogenetics answers the analogous question exactly. The four-point
condition~\cite{buneman} holds on every tree metric,
\begin{equation}\label{eq:fp}
  d(a,b) + d(c,e) \;\le\; \max\bigl(d(a,c)+d(b,e),\; d(a,e)+d(b,c)\bigr),
\end{equation}
and a \emph{quartet certificate} is four taxa whose six pairwise counts, over the
one denominator they share under complete deletion within the quadruple, refute
\eqref{eq:fp} as an integer comparison. Prior work in this programme built such
certificates, proved their relationship to split
decomposition~\cite{bandeltdress}, and emitted them for kernel
rechecking~\cite{paperf,paperg}.

The transfer is immediate once one notices the following.

\begin{definition}[Behavioural profile]\label{def:profile}
Run a candidate program against $T$ fixed probes and record one symbol per probe:
\texttt{A} correct, \texttt{C} finite but incorrect, \texttt{G} raised,
\texttt{T} non-finite, and \texttt{-} where the probe is inadmissible.
\end{definition}

That is a four-state alphabet with a gap symbol, which is the model a nucleotide
alignment assumes; the gap is handled by the same complete deletion that handles
a missing base. A $\text{programs} \times \text{probes}$ matrix is thus a valid
alignment and the machinery runs on it with no modification.

The null is better than convenient, it is correct: a population whose only
variation is mutation from single parents has a genealogy that is a tree
\emph{by construction}. So the question ``does this crossover operator
recombine?'' becomes ``does the crossover-enabled population's margin exceed the
floor of a mutation-only population at matched budget?'' --- with an exhibit a
kernel can recheck.

\section{The synthetic test, which appeared to confirm it}

Eight candidates, $1200$ probes, three seeds per cell, two arms at matched
budget: a mutation-only lineage, and the same lineage with crossovers at interior
breakpoints. Mutation rate swept, because a statistic of this kind dies at both
ends.

\begin{center}\scriptsize
\begin{tabular}{@{}rllrrcc@{}}
\toprule
rate & clonal violating & recomb.\ violating & max clon.\ margin & min rec.\ margin & margin sep. & null varies \\
\midrule
$0.0005$ & $[0,0,0]$ & $[0,0,0]$ & $0$ & $0$ & no & \textbf{no} \\
$0.0020$ & $[0,0,0]$ & $[21,25,27]$ & $0$ & $0.00342$ & yes & \textbf{no} \\
$0.0080$ & $[13,19,21]$ & $[37,53,45]$ & $0.00342$ & $0.00514$ & \textbf{yes} & \textbf{yes} \\
$0.0200$ & $[60,61,56]$ & $[59,59,65]$ & $0.00685$ & $0.01370$ & \textbf{yes} & \textbf{yes} \\
\bottomrule
\end{tabular}
\end{center}

Three readings, each a known lesson of this programme reappearing. Certificates
are found and all revalidate. The \emph{margin} is the statistic and the count is
not: at rate $0.02$ the count fails outright (clonal $56$--$61$ against
recombinant $59$--$65$) while the margin still separates twofold. And both
extremes are unusable for opposite reasons --- below $0.002$ the clonal null
produces no violating quadruple at all, so any positive observation is trivially
significant and must be reported as degenerate rather than as a result.

This is the experiment that should have been distrusted, and \S\ref{sec:diagnosis}
says why.

\section{The real test}\label{sec:real}

\paragraph{Representation.} Expression trees over a closed grammar: terminals
$x$ and small integer constants, operators $+$, $-$, $\times$ and a protected
division that \emph{raises} on a zero denominator rather than returning a
fallback, so a structural failure stays distinguishable from a wrong answer. The
grammar cannot express an import, an attribute access or a call, so individuals
are safe to evaluate directly and no sandbox is load-bearing.

\paragraph{Task and probes.} Target $(x^3 - 2x + 1)/(x^2+1)$, which makes
division genuinely necessary so the \texttt{G} state is reachable. $600$ probes
drawn uniformly on $[-2,2]$; tolerance $0.01$.

\paragraph{Arms.} Both spend the same number of variation events per generation
and cap program size, so neither event budget nor bloat can confound. The clonal
arm applies point mutations only --- a tree by construction. The recombinant arm
routes a fraction of its events through subtree exchange, which is a
recombination with a breakpoint; that fraction is swept over
$\{0.1, 0.2, 0.4, 0.6\}$.

\paragraph{Pre-registered before any run.}
\begin{itemize}
\item[\textbf{P1}] the clonal null is non-degenerate: it produces violating
quadruples, so something can be measured against it;
\item[\textbf{P2}] the floor converges: $25$ clonal runs do not raise the maximum
above the first $5$;
\item[\textbf{P3}] the recombinant margin exceeds the clonal floor, decided on
the integer count first and only then as a rate;
\item[\textbf{P4}] the effect is not an artifact of fitness: it survives
restricting both arms to comparable best fitness;
\item[\textbf{P5}] the effect is not an artifact of divergence: it survives
restricting the recombinant arm to runs whose \emph{realised behavioural spread}
lies inside the clonal arm's range.
\end{itemize}

\paragraph{Result.} P1 holds: $5$ of $25$ clonal runs carry certificates, with up
to $30$ violating quadruples. P2 holds exactly: the floor at $5$ runs and at $25$
runs are both $0.015000$, a growth factor of $1.00$.

\begin{center}\footnotesize
\begin{tabular}{@{}lcc@{}}
\toprule
pre-registered prediction & held & of \\
\midrule
P3 recombinant margin exceeds the clonal floor & $\mathbf{0}$ & $24$ \\
P4 \ldots restricted to comparable best fitness & $\mathbf{0}$ & $23$ \\
P5 \ldots restricted to comparable realised spread & $\mathbf{0}$ & $22$ \\
\bottomrule
\end{tabular}
\end{center}

$23$ of $24$ recombinant runs produce \emph{zero} violating quadruples, at every
crossover share. P5 matters most: $22$ of the $24$ were inside the clonal spread
range, so the failure is not a matching artifact. Crossover also bought no
fitness --- clonal median best error $0.5707$ against recombinant $0.6028$, with
the best individual overall identical in both arms at $0.4745$.

\paragraph{What actually drives the certificates.} Pooling all $49$ populations
and stratifying by realised behavioural spread:

\begin{center}\footnotesize
\begin{tabular}{@{}lrrrr@{}}
\toprule
spread & arm & runs & with certificates & mean violating \\
\midrule
$[0, 0.01)$ & clonal & $12$ & $0$ & $0.0$ \\
$[0, 0.01)$ & recombinant & $17$ & $0$ & $0.0$ \\
$[0.01, 0.05)$ & clonal & $9$ & $1$ & $1.2$ \\
$[0.01, 0.05)$ & recombinant & $6$ & $0$ & $0.0$ \\
$[0.15, 1.00)$ & clonal & $4$ & $\mathbf{4}$ & $21.8$ \\
$[0.15, 1.00)$ & recombinant & $1$ & $\mathbf{1}$ & $9.0$ \\
\bottomrule
\end{tabular}
\end{center}

Above a spread of $0.05$ every population carries a certificate in both arms;
below $0.01$ none of $29$ does in either. \textbf{Behavioural spread accounts for
certificate presence and arm membership adds nothing} --- and the clonal arm,
which cannot recombine, carries them at the same rate.

\section{Diagnosis: heritability is the load-bearing assumption}\label{sec:diagnosis}

The four-point condition detects non-tree-like structure \emph{in a distance
matrix}. It detects \emph{recombination} only under an additional assumption that
nucleotide data satisfies and behavioural profiles do not.

\begin{remark}[Why recombination is visible in DNA]
A nucleotide site is an independently heritable unit: a child's state at site $i$
is its parent's state at site $i$, modified rarely and locally. Recombination
produces \emph{mosaic ancestry along the sequence} --- a prefix from one parent
and a suffix from another --- so different sites carry different genealogies and
the resulting distances cannot be embedded in one tree. That conflict is what
\eqref{eq:fp} detects.
\end{remark}

A probe outcome has no such property. It is not inherited; it is \emph{computed}
from the whole program. A single point mutation near the root of an expression
tree changes the outcome at many probes at once, and two unrelated mutations in
two lineages routinely change overlapping probe sets. That manufactures exactly
the pattern of conflicting distances that \eqref{eq:fp} reports, with no
recombination anywhere --- which is precisely what the clonal arm's $4/4$ at high
spread shows.

\paragraph{The synthetic validation was circular.} Its profiles were produced by
mutating profile \emph{strings} position by position, which \emph{builds
site-wise heritability in by construction}. It therefore confirmed that the
method detects recombination in site-wise-heritable data --- which was never in
doubt --- while appearing to validate an application in which that property
fails. This is a non-discriminating check of the kind this programme has had to
label before: a test both hypotheses pass. The distinguishing experiment is the
one in \S\ref{sec:real}, where profiles are \emph{measured} from executed
programs instead of generated, and it reverses the conclusion.

\section{What survives}

\paragraph{The proof-carrying fitness ladder is untouched.} The second half of
the proposal replaced an empirical objective with one whose upper rungs are
Lean~4 obligations read through \texttt{\#print axioms}, so that no amount of
latency gain buys a rung and \texttt{sorry} cannot pass. Nothing here bears on
it, and the motivation stands: an un-kernel-checked gate is eventually found to
be forced, degenerate, or measuring the apparatus.

\paragraph{Three guards earned their place.} The programme's measurement rules
caught this rather than ratifying it: the demand that a null be able to vary
flagged the degenerate low-rate cells; the demand that a floor be shown converged
is what made $0/24$ trustworthy; and the demand to match on the \emph{realised}
nuisance rather than the nominal budget produced P5, without which the failure
would plausibly have been dismissed as a spread artifact.

\paragraph{A released dataset.} $49$ evolved program populations --- $25$ clonal
and $24$ recombinant across four crossover shares --- each with its behavioural
profile matrix in FASTA, violating-quadruple counts, exact margins, realised
spread, distinct-profile counts, best fitness and the best program rendered as an
expression. Useful to anyone asking whether a crossover operator contributes,
independently of the method refuted here.

\paragraph{A transferable negative.} Before carrying a phylogenetic method to a
non-biological population, ask what the method assumes about \emph{heritability
of characters}, not merely about the alphabet. Matching the data model --- four
states and a gap --- was necessary and nowhere near sufficient.

\section{Limitations}

One task, one grammar, one population size, and tournament selection with
elitism; subtree exchange is one crossover among several. A different
representation with genuinely modular, independently heritable components ---
separate functions recombined as units, say --- might restore the assumption
\S\ref{sec:diagnosis} identifies as missing, and that is the obvious next
experiment rather than a defence of this one. The negative result is about the
transfer as proposed and tested, not a theorem about all program populations.

The fitness comparison is also underpowered for its own question: $25$ against
$24$ runs on one task cannot settle whether crossover helps GP in general, and no
such claim is made.

\begin{thebibliography}{9}
\bibitem{buneman} Buneman P (1971/1974). The recovery of trees from measures of
dissimilarity. In \emph{Mathematics in the Archaeological and Historical
Sciences}, Edinburgh University Press, 387--395.
\bibitem{bandeltdress} Bandelt H-J, Dress AWM (1992). A canonical decomposition
theory for metrics on a finite set. \emph{Advances in Mathematics} 92(1),
47--105. \url{https://doi.org/10.1016/0001-8708(92)90061-O}
\bibitem{paperf} Callens X (2026). Proof-carrying detection of reticulate
evolution. Edition 8. Zenodo. \url{https://doi.org/10.5281/zenodo.23610116}
\bibitem{paperg} Callens X (2026). A proved lift that does not transfer:
split-level certificates of non-tree-likeness. Edition 1. Zenodo.
\url{https://doi.org/10.5281/zenodo.23604793}
\end{thebibliography}

\end{document}
